\section{Butterworth Highpass and AC-Coupled In-Amp}

\subsection{Butterworth Highpass}

\subsubsection{Second-order Butterworth highpass stage}
\[
H_{\mathrm{HP}}(s)=K\frac{s^2}{s^2+a\omega_c s+\omega_c^2}.
\]
The highpass numerator places two zeros at the origin. The denominator uses the same Butterworth damping coefficients as the lowpass prototype.

\subsubsection{Lowpass-to-highpass Butterworth transformation}
\[
H_{\mathrm{HP}}(s)=H_{\mathrm{LP,n}}\!\left(\frac{\omega_c}{s}\right).
\]
This substitution maps the normalized Butterworth lowpass prototype to a highpass cutoff at \(\omega_c\).

\subsubsection{Highpass cutoff conversion}
\[
\omega_c=2\pi f_c.
\]
Filter design often specifies cutoff frequency in hertz, while transfer functions use angular frequency in rad/s.

\subsection{AC Coupling}

\subsubsection{Single-pole AC-coupling highpass}
\[
H_{\mathrm{AC}}(s)=\frac{sRC}{1+sRC}
=\frac{s}{s+\omega_c},\qquad \omega_c=\frac{1}{RC}.
\]
AC coupling rejects electrode offsets and slow baseline drift while passing higher-frequency signal components.

\subsubsection{Cascaded AC-coupling stages}
\[
H_{\mathrm{total}}(s)=\prod_{k=1}^{M}\frac{s}{s+\omega_{c,k}}.
\]
Multiple coupling stages increase low-frequency attenuation and add phase shift near each cutoff.

\subsection{Instrumentation Amplifier}

\subsubsection{Differential and common-mode signals}
\[
v_d=v_2-v_1,\qquad
v_{\mathrm{cm}}=\frac{v_1+v_2}{2}.
\]
Biopotential front ends separate the desired differential signal from common-mode interference.

\subsubsection{Differential gain and common-mode gain}
\[
v_o=A_d v_d + A_{\mathrm{cm}}v_{\mathrm{cm}}.
\]
An ideal instrumentation amplifier has large differential gain and zero common-mode gain.

\subsubsection{Common-mode rejection ratio}
\[
\mathrm{CMRR}=\left|\frac{A_d}{A_{\mathrm{cm}}}\right|,
\qquad
\mathrm{CMRR}_{\mathrm{dB}}=20\log_{10}\left|\frac{A_d}{A_{\mathrm{cm}}}\right|.
\]
CMRR quantifies rejection of common-mode signals relative to differential signals.

\subsubsection{Three-op-amp instrumentation amplifier gain}
\[
A_d=\left(1+\frac{2R}{R_G}\right)\frac{R_3}{R_2}.
\]
The gain is set by the input-stage gain resistor \(R_G\) and the output difference-amplifier resistor ratio.

\subsubsection{Resistor mismatch and common-mode leakage}
\[
\frac{R_3}{R_2}\neq \frac{R_5}{R_4}\quad \Rightarrow\quad A_{\mathrm{cm}}\neq 0.
\]
Balanced resistor ratios are required for common-mode cancellation. Mismatch degrades CMRR.
